\documentclass[12pt]{article}
\usepackage[T1]{fontenc}
\usepackage[utf8]{inputenc}
\usepackage{lmodern}
\usepackage{textcomp}
\usepackage{enumitem}
\usepackage{geometry}
\usepackage{graphicx}
\usepackage{amsmath, amssymb}
\geometry{margin=1in}
\begin{document}
\begin{center}
Economics - Graduate Level Assessment 21 \
Total Marks: 40 \
Answer all questions.
\end{center}

\section*{Multiple Choice (10 marks)}
\begin{enumerate}[label=\arabic*.]
\item Which of the following would shift the aggregate demand curve to the right?
\begin{enumerate}[label=(\alph*)]
    \item An increase in unemployment rates
    \item A decrease in government spending
    \item A decrease in the supply of money
    \item An increase in savings
    \item An increase in interest rates
\end{enumerate}
\item Does the principle of comparative advantage apply only to international trade?
\begin{enumerate}[label=(\alph*)]
    \item No, it is a principle that applies only to consumer choice in the market.
    \item No, it is only applicable to the exchange of services, not goods.
    \item Yes, it specifically applies to trade agreements between neighboring countries.
    \item No, it applies wherever productive units have different relative efficiencies.
    \item Yes, it only applies to international trade.
\end{enumerate}
\item If Matt's total utility from consuming bratwurst increased at a constant rate, no matter how many bratwurst Matt consumed, what would Matt's demand curve for bratwurst look like?
\begin{enumerate}[label=(\alph*)]
    \item Linearly increasing
    \item Inverted U-shape
    \item Horizontal
    \item Upward sloping
    \item Vertical
\end{enumerate}
\item The aggregate demand curve is
\begin{enumerate}[label=(\alph*)]
    \item not found by adding product demand curves horizontally or vertically
    \item a vertical summation of market demand curves
    \item found by adding product demand curves both horizontally and vertically
    \item a horizontal subtraction of firm demand curves
    \item a simple aggregation of demand curves for individual goods
\end{enumerate}
\item What school of economists defended Say's Law? Which pro-minent economist disputed this theory?
\begin{enumerate}[label=(\alph*)]
    \item Chicago school, Stiglitz
    \item Marxist school, Minsky
    \item Behavioral school, Akerlof
    \item Neoclassical school, Friedman
    \item Neoclassical school, Coase
\end{enumerate}
\end{enumerate}

\section*{True / False (10 marks)}
\textit{Answer True or False}
\begin{enumerate}[label=\arabic*.]
\setcounter{enumi}{5}
\item True or False: Recovery from a recession, such as the American recession of 1974-75, is often classified as economic reformation.
\item True or False: The economic growth from 1990 to 2000, based on real GDP per capita, was 75 percent.
\item True or False: A perfectly competitive firm will continue to operate as long as the price exceeds average total cost, regardless of variable costs.
\item True or False: Classical analysis suggests that increasing government spending is the primary means to restore full employment during a recession.
\item True or False: Marginal cost always intersects average variable cost at the profit-maximizing quantity of output.
\end{enumerate}

\section*{Long Form Response (20 marks)}
\textit{Answer each question with a clear and complete explanation.}
\begin{enumerate}[label=\arabic*.]
\setcounter{enumi}{10}
\item Discuss the role of money as a unit of account, using \$1.25 as the price of a 20-ounce bottle of pop to illustrate its function in facilitating economic transactions.
\item Discuss the conditions under which marginal cost equals average variable cost and average total cost, and analyze the implications of these intersections on production decisions.
\end{enumerate}

\vfill
\noindent\textit{End of Paper}
\end{document}